\section{Conclusion}\label{section:Conc_Recc}
%%insert recs intro
This article describes the opportunities and challenges of spiking neural networks (SNN) for closed-loop control, trained with actor-critic reinforcement learning. While the convergence of SNN during training was slow and noisy, the trained agents were able to control the CartPole. Evaluating the trained agents using NeuroBench highlights the use cases for ANN as well as SNN. Currently, the consistency and performance of ANN allow for a more reliable and performing controller in simulation. However, two aspects of the spiking agents offer interesting applications for systems deployed in real life. First, the computational resources required for SNN showed significantly lower than their non-spiking counterpart. These resources were further reduced by applying a simple pruning algorithm, leading to a 25x reduction in operations required per control action. This allows the agent to be run on extremely low power devices, or alternatively run the model with a much higher frequency if necessary. Next, the presence of noise in the observation of the agent severely affected the ANN models. At relatively low gains of the noise inserted, the SNN trained with temporal dynamics ($\beta>0$) outperformed the ANN. Therefore, in real systems, the expected noise should be analyzed before the deciding on spiking or non-spiking controllers.

Further research is needed to improve the training efficiency of SNN using RL. Techniques such as reward-modulated STDP and eProp show promise in this regard. Additionally, the noise characteristics showed promising results for real-world applications. Future work could use the presented methods to deploy a controller in the real world and analyze their performance. Overall, the exploration of spiking agents in reinforcement learning is an interesting avenue for closed-loop control on deployed systems and warrants further investigation.